\documentclass{article}

\usepackage[utf8]{inputenc}     % Soporte para tildes y eñes
\usepackage[T1]{fontenc}
\usepackage[spanish]{babel}     % Para usar títulos en español si querés
\usepackage{amsmath, amssymb}   % Símbolos matemáticos
\usepackage{xcolor}             % Colores para código
\usepackage{listings}           % Código fuente
\usepackage{geometry}           % Márgenes
\usepackage{array}


\geometry{margin=2.5cm}

% Configuración para mostrar C++ lindo
\lstset{
  language=C++,
  basicstyle=\ttfamily\small,
  keywordstyle=\color{blue}\bfseries,
  commentstyle=\color{gray},
  stringstyle=\color{red},
  numbers=left,
  numberstyle=\tiny\color{gray},
  stepnumber=1,
  breaklines=true,
  tabsize=2,
  showstringspaces=false,
  captionpos=b
}

\title{Matemática}
\author{Francisco Carthery}
\date{July 30, 2025}

\begin{document}

\maketitle

\section{Teoría de números}

\subsection{Factorizacion en primos}
\[
n = p_2^{\alpha_1} p_2^{\alpha_1} \cdots p_i^{\alpha_i}
\]


\subsection{Numero de factores}
\[
\tau(n) = \prod_{i=1}^{k} (\alpha_i + 1)
\]

\subsection{Suma de los factores}
\[
\sigma(n) = \prod_{i=1}^{k} (1 + p_i + \dots + p_i^{\alpha_i}) = \prod_{i=1}^{k} \frac{p_i^{\alpha_i + 1} - 1}{p_i - 1}
\]

\subsection{Producto de los factores}
\[
\prod d = \prod_{i = 1}^{k} p_i^{\frac{\alpha_i (\alpha_i + 1)}{2} \frac{\tau (n)}{\alpha_i + 1}} = N^{\frac{\tau(n)}{2}}
\]

\subsection{Cantidad de primos hasta n}
\[
\pi(n) \approx \frac{n}{\ln(n)}
\]

\subsection{Funcion phi de Euler}
Nos dice la cantidad de enteros entre 1 \dots n que son coprimos con n
\[
\phi(n) = \prod_{i=1}^{k} p_i^{\alpha_i - 1} (p_i - 1)
\]

\subsection{Teorema de Euler}
Para todo par de \textbf{coprimos} positivos \textit{x} y \textit{m}:
\[
x^{\phi(m)} mod \: \textit{m} = 1
\]


\begin{center}
\begin{tabular}{>{$}l<{$}|*{15}{c}}
\multicolumn{1}{l}{$n$}                     \\\cline{1-1}
0  & 1    &      &      &      &      &      &      &      &      &      &      &      &      &      &      \\
1  & 1    & 1    &      &      &      &      &      &      &      &      &      &      &      &      &      \\
2  & 1    & 2    & 1    &      &      &      &      &      &      &      &      &      &      &      &      \\
3  & 1    & 3    & 3    & 1    &      &      &      &      &      &      &      &      &      &      &      \\
4  & 1    & 4    & 6    & 4    & 1    &      &      &      &      &      &      &      &      &      &      \\
5  & 1    & 5    & 10   & 10   & 5    & 1    &      &      &      &      &      &      &      &      &      \\
6  & 1    & 6    & 15   & 20   & 15   & 6    & 1    &      &      &      &      &      &      &      &      \\
7  & 1    & 7    & 21   & 35   & 35   & 21   & 7    & 1    &      &      &      &      &      &      &      \\
8  & 1    & 8    & 28   & 56   & 70   & 56   & 28   & 8    & 1    &      &      &      &      &      &      \\
9  & 1    & 9    & 36   & 84   & 126  & 126  & 84   & 36   & 9    & 1    &      &      &      &      &      \\
10 & 1    & 10   & 45   & 120  & 210  & 252  & 210  & 120  & 45   & 10   & 1    &      &      &      &      \\
11 & 1    & 11   & 55   & 165  & 330  & 462  & 462  & 330  & 165  & 55   & 11   & 1    &      &      &      \\
12 & 1    & 12   & 66   & 220  & 495  & 792  & 924  & 792  & 495  & 220  & 66   & 12   & 1    &      &      \\
13 & 1    & 13   & 78   & 286  & 715  & 1287 & 1716 & 1716 & 1287 & 715  & 286  & 78   & 13   & 1    &      \\
14 & 1    & 14   & 91   & 364  & 1001 & 2002 & 3003 & 3432 & 3003 & 2002 & 1001 & 364  & 91   & 14   & 1    \\
\hline
\multicolumn{1}{l}{} 
   & 0 & 1 & 2 & 3 & 4 & 5 & 6 & 7 & 8 & 9 & 10 & 11 & 12 & 13 & 14 \\\cline{2-16}
\multicolumn{1}{l}{} &\multicolumn{15}{c}{$k$}
\end{tabular}
\end{center}






\end{document}
